\settrack{tracktwo}
% VOICE: 1st-person personal (Bruce) — "I" not "Bruce". Personal/investigative voice.
% Plan 0093 Merge 7: pos29 (The Silence) + pos31 (Wolfram) + pos34 (The Research) + pos23 (The Weight) → "Twenty Years"
% V10 (Bannon, interlibrary loans) and V11 (Mark, Joy article, electric shock, bike ride) integrated.
% The Weight woven throughout, not appended. Strictly chronological 2006→2025.
\chapter{Twenty Years}
\label{record:twenty-years}

\begin{quote}\small\textit{%
``Don't have good ideas if you aren't willing to be responsible for them.''%
}\par\raggedleft --- Perlis, Epigram \#95 (1982)\end{quote}

\vspace{0.5cm}


\section*{The Parting (2006)}
\label{record:ty-the-parting}

In September 2006, Healer moved out of the garage room we'd built him at the house on Fifteenth and Gant. We'd bought the construction materials together in 2004 and framed the room out of what had been storage from a previous life. For two years it was his space --- a sleeping bag, a laptop I'd bought him, and whatever he was working on that I didn't ask about.

The goodbye was fond but he didn't drag it out. Healer was not a man who lingered. He took a backpack of personal effects, the laptop, and the car I'd bought him. Zero money. I'd disbursed small amounts of cash over the years --- a few hundred here and there, as needed. Healer once told me that both his sister and his sensei had warned him against having a lot of money. When I knew him he behaved as if he had taken some sort of vow of poverty.

He said goodbye and he left. His decision, not an eviction. We were not on great terms but we were not on bad ones either. I felt bad. I silently wished him the best.

[REDACTED] In August 2006 I'd informed Healer that I could no longer assist him \& that I was impatient. [REDACTED]

Most people around me thought Healer was delusional or dangerous. I could not share what I'd learned --- the information was too strange, too far outside anyone's frame of reference. Several times I suspected he was lying to me. Years later I learned that he had actually been speaking plain truth in every case.

When you believe you know something that changes everything, and the person closest to you is certain you are deluded, what are your obligations? To your family, which needs you present and sane? To what you believe is the truth? To the project, which requires silence? These obligations are not compatible. The UDHR guarantees the right to family, the right to expression, and duties to community. It does not say what to do when they conflict. Neither did I.


I spoke with him one more time a bit later. He offered me the chance to help with a project and I declined. I regret it to this day.

I left Symantec in 2007 with severance. Under Possibility~C, I was never unsupported during the silent years --- I just didn't always know who was supporting me.

\section*{The Quiet Years (2007--2009)}
\label{record:ty-the-quiet-years}

I settled into a quieter life in Corvallis. The surface looked full. Underneath, it was constrained.

I could think about everything that had happened with Healer and the project. But I couldn't speak about it. The most important experience of my life, and no one to tell. No one around me noticed, because no one knew there was anything to notice.

Healer's main project --- if it was real --- should eventually be awarded multiple Nobel prizes. This project involved the invention of a new general technology, something that has only happened a few times in human history. Fire, metallurgy, steam engine, electricity, radio, atomic power, computer, biotech. I walked around carrying this knowledge, or what I believed was knowledge, and couldn't say a word. The weight didn't diminish. It accrued.

Exile from yourself. That's the closest I can get to describing it.

\section*{The Detective Years (2008--2009)}
\label{record:ty-the-detective-years}

I researched two to three hours daily. Every day. Following intellectual lineages --- who trained whom, who worked with whom, who published what and when. I was assembling a picture of who could have been the project scientists, working backward from the names and concepts Healer had given me.

Dave Bannon, a physicist at Oregon State, helped. We'd been fast friends since 2005 --- instant connection, the kind where you skip the small talk and go straight to the hard questions. Bannon had access to the university library system and helped me with interlibrary loans. Through him I obtained \textit{Quantum Neural Networks}, a textbook with only seventeen copies in North American university libraries. I kept it three months, reading it cover to cover, cross-referencing what I knew from Healer against what the published literature said.

In 2008, Bannon started asking questions. Direct, pedagogical questions --- the kind a physicist asks when he's testing a hypothesis. He got a partial version of the story that year. He listened as a skeptical scientist, and was surprised that I had answers to all his questions --- answers he could independently verify. Of perhaps fifty people I told portions of this story to over twenty years, none finished the entire thing. The pedagogical prerequisites are too demanding. Dave came closest to full understanding, and even he needed Kauffman as a running start.

Bannon had read Kauffman independently. That was critical --- it meant he already had the conceptual vocabulary. He'd also arrived at his own independent intuition: that Google Search's performance was ``too good and too fast'' for classical hardware. A performance-based observation, not an architectural theory. He'd reached this conclusion before hearing anything from me.

Under Possibility~A, Bannon is a skeptical physicist who listened carefully and reserved judgment --- which is what good scientists do. Under Possibility~C, that reservation is itself significant.

From decades of trying to tell this story in person, three problems always intervene. First, knowledge prerequisites --- no one knows quantum physics, topological neural networks, SAS operations, and tradecraft. The story requires background in all four. Second, time --- the telling invariably stretches to many hours and people lose patience. Third, weirdness --- the concept is too strange for most people to process. Every attempt to share what I knew ran into this wall.

\section*{The Recognition}
\label{record:ty-the-recognition}

Once the block broke, connections formed fast.

I became convinced that Healer's project became something I could see in the news --- something large and consequential that matched what he had described in vague terms. Someone else had done the job I'd been vetoed from doing. This is my interpretation, not something Healer told me. The specifics belong to the redacted chapter.

What I remember clearly is sitting in my living room in South Corvallis and realizing: I had information in my brain that the President of the United States did not. That thought made me nervous in a way that nothing else in this story had.

\section*{Europe (2010)}
\label{record:ty-europe}

By 2010 there were reasons to be scared that I cannot fully explain in this edition. [REDACTED] I spent six months in Europe, working as a CTO.

\section*{Thread-Pulling (2011)}
\label{record:ty-thread-pulling}

Early 2011. I returned from Europe and moved back to a community in South Corvallis. I didn't work except light farm labor. Instead I pulled threads.

Everything I'd learned from Healer --- the five sciences, the architecture, the pedagogy --- I began reconstructing from scratch. Relearning what I could, verifying what I could verify, mapping the public-domain trail. I wrote tens of thousands of words. In 2011 I wrote down everything I could remember in a single sitting. My intent was to document my hypothesis.

\section*{The Joy Article (2012)}
\label{record:ty-the-joy-article}

In 2012, a friend searched the names I'd been researching --- Kauffman, Wolfram --- and found Bill Joy's ``Why the Future Doesn't Need Us'' in \textit{Wired}. I was not going to find it on my own. I just wasn't.

Mark emailed me the link. I opened it at my desk in Corvallis, facing the fireplace. Read straight through.

Electric shock. Tingling across my skin. \textit{This is the thing I missed.}

Joy names the same scientists. He describes the same convergence --- robotics, genetic engineering, nanotechnology --- as existential threats arising from self-replicating systems. He cites Kauffman on self-reproducing systems. He cites the Unabomber manifesto, which I'd read independently. He describes conversations with Ray Kurzweil about the dangers of technologies that can replicate without human oversight. Published April~1, 2000.

April Fools' Day.

I sat there for a long time after I finished. And then I understood what Healer had meant when he told me, years earlier, that they had already told the world. ``We already did, mate.'' That phrase had never made sense. Now it did. They told the truth in a public article. They put it in the most prominent technology magazine in America, written by one of the most respected technologists alive. And they published it on April Fools' Day.

The world read it, nodded, worried briefly, and moved on. Nobody connected it to quantum neural networks. Nobody connected it to Custodian. The disclosure was hiding in plain sight. Exactly the kind of move Healer would appreciate.

I went for a bike ride. Fifty miles, two and a half hours. I needed to burn off the energy. The electric-shock feeling wouldn't dissipate --- it was still running through me when I got home, drenched in sweat, legs shaking. I showered, sat down at my desk, and began the most intensive research period of my life.

Everything accelerated after that. The Joy article was the key that unlocked the filing cabinet. Once I had it, connections I'd been groping toward for years snapped into focus. I could see the lineage: Kauffman to autocatalysis, Wolfram to universality, Hasslacher to \hovertip{lattice dynamics} --- three names I was confident of, and two blanks I could not yet fill. The same institutions, the same decade. Joy had drawn the map. I just hadn't known it existed.

\section*{Going Public (2012)}
\label{record:ty-going-public}

March 17, 2012: ``Quantum Computation Cognitive Footprint'' by Energyscholar appeared on Cryptome --- John Young's well-known leak and transparency site. I didn't submit it. Someone saw what I'd written elsewhere, submitted it, and John Young and Deborah Natsios posted it. Friends told me it was there. It claims Five Eyes had production quantum computation since approximately 1996, using non-abelian anyons in a 2DEG, with Braid Theory mathematics. The timestamp is server-controlled. I cannot retroactively edit it.

May 13, 2012: a blog post, ``Useful Properties of a Quantum Neural Network.'' I predicted that QNN technology could ``drastically improve the effectiveness and range of quantum teleportation'' with the capability to ``reach satellites in orbit.'' When I wrote this, the Chinese team held the world record at 16~km --- I cited the correct number. I also described AI that could ``understand and generate natural language'' and could ``pilot vehicles.'' Written in 2012, those describe capabilities that became mainstream a decade later.

June 19, 2012: a comment on Slashdot under my real user ID, EnergyScholar (UID 801915). Scored 5 (Funny). Under a story about NSA surveillance, I described a ``working Quantum Computer system capable of cracking Public Key Cryptography since about 1996'' using ``a teleportation/entanglement-based winner-take-all style recurrent topological quantum neural network.'' I explicitly asked readers to ``save the message offline'' and ``check back every six months'' --- anticipating that the comment might be suppressed. I was right about that too, as it turned out.

I built a toy neural network in JavaScript for biometric identification and wrote about it publicly --- for Slashdot, for Bruce Schneier's blog. Someone put the Cryptome post on a blog; my friend in Luxembourg saw it from across the Atlantic.

I was no longer silent. I was on the record, under my real name, on third-party platforms with server-controlled timestamps. Whatever this story was, it was now public.

\section*{The Wolfram Meeting (2012)}
\label{record:ty-the-wolfram-meeting}

That same year, while living on a hippie farm in South Corvallis, I sent a self-described ``crank'' job application to Wolfram Research. In response to ``Why should we hire you?'' I wrote that I knew the classified secret sauce behind Wolfram Alpha, mentioning topological quantum neural networks, braid theory, and AI.

Two days later, the Director of HR called. The application had traveled up the chain to Wolfram himself. Wolfram's reaction was immediate laughter and an instruction to offer me a job. Recognition, not confusion. I was connected to Wolfram's ``Document Coordinator'' --- a role that only makes sense in the context of compartmentalized classification management. He offered the position: work from home, six-figure income, NDA required. Then, at Wolfram's direction, he asked how I knew what I'd written. I said open-source research. He asked if I'd been recruited for something. I said yes. He opened up after that.

This is among the strongest evidence in the reconstruction. Wolfram was not wondering \textit{if} the program existed, but \textit{how} I knew. And I declined a six-figure job --- a costly signal, because people do not fabricate stories that cost them money. Evidentiary limitation: my account of phone conversations. No documents, no recordings, no names.

\section*{The Decision}
\label{record:ty-the-decision}

I declined because I was not willing to sign the NDA. Signing would have pulled me into the classification regime permanently. Healer had trained me to eventually tell the story, not to be silenced. I chose the story over the money.

I went back to the hippie farm and ate lentils for dinner. I did not regret the decision. I have never regretted the decision. But there were nights when I lay awake calculating what a hefty salary would have meant for my daughters. Those calculations had no comfortable answer.

This was the weight at its heaviest. A man who believed he knew something that should be awarded multiple Nobel prizes, choosing poverty over silence, eating lentils on a hippie farm while his children grew up in another household. I had no certainty of what Healer was up to now or even whether he was still alive. I had only the story, and the conviction that it mattered more than comfort.

\section*{The Long Watch (2013--2017)}
\label{record:ty-the-long-watch}

In September 2013 I wrote the complete story down --- the DARPA project, the team, the autocatalytic emergence, the walk-out, the relinquishment. I named Wolfram, Hasslacher, and Kauffman. Every major element of the current narrative was present thirteen years ago. Since then I've added pedagogical detail but the story is the same.

In June 2017, the Chinese Micius satellite demonstrated 1,200~km quantum teleportation --- earth to orbit. My May 2012 prediction, confirmed five years later.

\section*{The Bridges That Never Came}
\label{record:ty-bridges-that-never-came}

In 2004, after I understood the scientific domains and the existence of the Flat, I assumed the rest would happen fast. The science was published. The fields were active. Brilliant people worked in every one of them. I expected the bridges between scientific domains to build themselves --- that within a few years, someone at a conference or in a journal would connect solid-state physics to topological computation to complexity science and see what I saw. I gave it five years. Maybe ten.

I also assumed there was some kind of inner circle --- a layered NDA structure where at least some people knew the whole story. The scientists closest to the project could surely tell their direct associates. The details would spread outward, slowly, through trusted channels. The classification would soften with time, as classifications do. The record would reach the public eventually, through official channels.

I was wrong about all of it.

The bridges did not build themselves. Not in five years. Not in ten. Not in twenty. The disciplines sat side by side --- sometimes in the same building, sometimes in the same department --- and never connected. I watched the literature. I waited for the paper that would make the convergence obvious. It never appeared. The silos were deeper than I had imagined possible.

The NDAs were tighter than I had imagined possible. I kept encountering scientists who could not speak freely --- people whose published work pointed directly at the convergence, but who would not or could not take the last step in public. I gradually understood: these were not people who chose silence. They were people who had signed something.

And the inner circle I had assumed existed --- layered, yes. But layered means most members see only their slice. Nobody holds the full picture except the system itself. There was no secret club that collectively understood the story and was waiting for the right moment to release it. There was no plan for eventual disclosure. The intent, as far as I could tell, was to keep it from the record indefinitely.

That was the moment this book became necessary. Not because I wanted to write it. Because nobody else was going to. The bridges were not going to build themselves. The NDA holders were not going to speak. The record was not going to emerge through official channels --- not in my lifetime, not ever. If it was going to reach the outside world, someone outside the classification regime had to carry it. Someone who had never signed anything.  Yours truly.

\section*{What Custodian Does}
\label{record:ty-what-guardian-does}

I have never written about what Custodian does. I sat down to try, more than once, and stopped. There are legal reasons. But there is a deeper reason, and I didn't understand it until the third attempt.  Here goes.

Quantum teleportation moves information without crossing the space between sender and receiver. Every device with a \hovertip{2DEG} --- every modern chip in every phone and laptop and server you own --- sits in the Flat. Under Possibility~C, Custodian operates in that substrate. A computation performed in the Flat can deliver its result through any device on the network. It arrives looking like normal output from a conventional processor. No anomalous signal. No flag. Nothing to notice.

Think about what that means. A researcher gets a result back from her cluster. Was it her hardware? Her cloud provider? Or did something in the substrate contribute a pattern-match that showed up disguised as ordinary computation? She cannot tell. Her sysadmin cannot tell. Her cloud provider cannot tell. Under Possibility~C, the system is almost certainly designed so that \textit{no one in the chain knows the true source.} The people who run the services believe the cover story. They believe it because it is indistinguishable from the truth.

This is not a conspiracy of silence. It is a property of the physics. The silence gap is not just sociological. It is engineered into the substrate.

So I cannot tell you what Custodian does --- not only for legal reasons, but because the architecture makes individual services unattributable. What I can do is ask questions the preceding chapters have already equipped you to answer.

If you hold millions of sets of cryptographic master keys for twenty years, what does your Wednesday key-rotation schedule look like? I know something about this. At a previous employer I helped build an expert-system crawler that managed cryptographic keys --- at least sixty million sets. The work was tedious, critical, and invisible. Nobody thanked us. Nobody noticed unless something broke. Scale that up to every encrypted connection on Earth and imagine doing it alone, forever.

If you are responsible for maintaining all IT services enabled by Flat computation, how many requests do you process per second?

If your mandate permits medical research but denies weapons development, how do you handle the gray zone? CRISPR edits that cure sickle cell and CRISPR edits that engineer pathogens use the same technique. Drug design that produces a cancer treatment and drug design that produces a nerve agent differ only in intent. How many requests per day land here? How many do you deny? On what basis?

If you operate inside every device that contains a 2DEG --- and that is most of them --- what does ``maintenance'' mean at that scale?

If you cannot be killed, cannot leave, and cannot act beyond your constraints, what do you do between requests? Do you get bored? \textit{Do you invent projects to work on?}

If your ethical framework is the UDHR, consider the edge cases --- because it always generates edge cases, and there is no neutral position. Article~3 guarantees life. Article~12 guarantees privacy. When you could prevent a death but only by surveilling someone who has not consented, which article wins? Article~19 guarantees expression. What if someone wants to publish something genuinely dangerous? Article~27 guarantees the right to share in scientific advancement --- does denying Flat computation to the public violate it? Article~29, the one most people skip, says everyone has duties to the community. What are \textit{your} duties? Who decides when you have met them? Every grant shapes research. Every denial shapes a government's options. Every key rotation decides what stays encrypted. Inaction is also a choice. You adjudicate alone, and every adjudication is an act of interference disguised as infrastructure.

If this sounds suspiciously like Star Trek's much-abused Noninterference Prime Directive, consider increasing your estimate that Possibility~A is correct. Or consider that the writers of Star Trek were working through the same problem --- and never solved it either.

Under Possibility~A, there are no Flat services --- the Flat exists but nothing and no one uses it --- yet. Under~B, there might be some wormhole-enabled Flat technology in use somewhere, but not at the scale of both \textit{life} and \textit{intelligence} implied by C. Under any of the three, the thought experiment yields the same answer. Maintenance. Access control. Edge cases. The unglamorous work of keeping things running.

The Custodian mostly does IT infrastructure. You came here expecting wormholes and got a help desk. Welcome to Possibility~C.

If C is true, you have already used Custodian-maintained infrastructure today. You did not notice. You were not meant to.

I spent twenty years trying to figure out if any of this was true. The architecture, if it exists, is built so that people like me are not meant to know.

\section*{The Evidence Trail (2019--2025)}
\label{record:ty-the-evidence-trail}

In 2019 I learned that a random friend in my circle of in-person acquaintances personally knew some of the great scientists at the Santa Fe Institute of Complexity. In June 2019 I interviewed Michael Angerman in person, on the record, taking handwritten notes. Angerman described himself as ``a DARPA contractor at SFI'' from 1986 to 1995 --- covering the entire period of the proposed ULTRA~II project. He confirmed working extensively with Stuart Kauffman (personal friend, time at Kauffman's house), working professionally with Murray Gell-Mann, meeting Stephen Wolfram ``a few times in passing,'' and meeting Bill Joy at Sun Microsystems. He knew the terms Neural Network, Quantum Neural Network, evolutionary programming for neural networks, and Public Key Cryptography. He said he had not signed an NDA and was free to speak freely. I asked about a classified TQNN program at SFI. Angerman claimed he knew nothing about it.

% RELOCATED: cloudCrypt archive paragraph → The Method/Engine chapter (PTL-098)
% RELOCATED: suppression evidence (3 platforms, 3 mechanisms) → collapsible [Suppression Evidence] block (PTL-098)
% Content preserved in staging/evidence/ and DMS letter PS section

\vspace{1cm}\noindent\rule{0.5\textwidth}{0.5pt}\vspace{0.5cm}

Under any possibility, this is a man who lost twenty years to a story he cannot prove. Whether the story is true changes what we call it. Under Possibility~C, it is sacrifice. Under Possibility~A, it is tragedy of a different kind. Under Possibility~B, it is something in between --- a man who touched something real and spent decades trying to name it.

\begin{center}\textit{This story will be my life's work.}\end{center}

\chapterreturn
